\documentclass[12pt, twoside]{article}
%\documentclass[12pt, twoside]{article}
\usepackage[letterpaper, margin=1in, headsep=0.2in]{geometry}
\setlength{\headheight}{0.6in}
%\usepackage[english]{babel}
\usepackage[utf8]{inputenc}
\usepackage{microtype}
\usepackage{amsmath}
\usepackage{amssymb}
%\usepackage{amsfonts}
\usepackage[nomessages]{fp} %\FPeval{\var-name}{2*sin(pi/6)}
\usepackage{siunitx} %units in math. eg 20\milli\meter
\usepackage{yhmath} % for arcs, overparenth command
\usepackage{tikz} %graphics
\usetikzlibrary{quotes, angles, arrows, arrows.meta}
\usepackage{pgfplots}
\usepgfplotslibrary{fillbetween}
\pgfplotsset{compat=1.16}
\usepackage{graphicx} %consider setting \graphicspath{{images/}}
\usepackage{parskip} %no paragraph indent
\usepackage{enumitem}
\usepackage{multicol}
\usepackage{venndiagram}

\usepackage{fancyhdr}
\pagestyle{fancy}
\fancyhf{}
\renewcommand{\headrulewidth}{0pt} % disable the underline of the header
\raggedbottom
\hfuzz=2mm %suppresses overfull box warnings

\usepackage{hyperref}

\fancyhead[LO]{La Scuola d'Italia / Dr.\ Huson / IB Math 12 \\* 7 October 2026}
\fancyhead[RO]{\fbox{\begin{minipage}{6cm}\footnotesize Nickname:\\[0.5cm]\end{minipage}}}
\fancyhead[LE]{\fbox{\begin{minipage}{6cm}\footnotesize Nickname:\\[0.5cm]\end{minipage}}}
\fancyfoot[C]{\thepage}

\begin{document}
\subsubsection*{1.9 Classwork: Bivariate data and linear regression}

\emph{Use your GDC. Give answers exactly, or correct to three significant
figures. Write down the values you put into the calculator.}

\begin{enumerate}[itemsep=0.15cm]

%--- Problem 1: Pearson's r, describing correlation, y-on-x regression line, interpreting the gradient, interpolation, extrapolation [9 marks] ---
%--- Source: Original (freshly authored, AA SL 4.4), GDC Linear Regression (mx+b) from two lists ---
\item[1.]

A gelateria in Rome recorded the maximum temperature, $x$ $^{\circ}$C, and
the number of gelati sold, $y$, on eight days in June.

\begin{center}
\setlength{\tabcolsep}{7pt}\renewcommand{\arraystretch}{1.2}
\begin{tabular}{|l|c|c|c|c|c|c|c|c|}
\hline
Temperature, $x$ ($^{\circ}$C) & 18 & 20 & 21 & 23 & 25 & 28 & 30 & 32 \\
\hline
Gelati sold, $y$ & 150 & 158 & 181 & 185 & 214 & 226 & 262 & 258 \\
\hline
\end{tabular}
\end{center}

\begin{enumerate}[itemsep=0.1cm]
  \item Find Pearson's product-moment correlation coefficient, $r$.
  \item Describe the correlation between $x$ and $y$ (direction: positive
  or negative; strength: strong, moderate, weak or none).
  \item Write down the equation of the regression line of $y$ on $x$, in
  the form $y = ax + b$.
  \item Interpret the meaning of the value of $a$ in this context.
  \item Use your regression line to estimate the number of gelati sold on
  a day when the maximum temperature is $26\,^{\circ}$C.
  \item Explain why the regression line should not be used to estimate the
  number of gelati sold on a day in January when the maximum temperature
  is $10\,^{\circ}$C.
\end{enumerate}

\begin{tikzpicture}
  \draw (-0.6,0) rectangle (14.6,10.5);
  \foreach \y in {1,2,...,10} \draw[dotted] (0,\y)--(14.1,\y);
\end{tikzpicture}

\newpage

%--- Problem 2: y-on-x line and interpreting b; mean point; x-on-y line; choosing the appropriate line to estimate x from y [9 marks] ---
%--- Source: Original (freshly authored, AA SL 4.4 and 4.10), GDC Linear Regression with the lists swapped ---
\item[2.]

A student asked eight classmates how many hours a day they spend on their
phone, $x$, and how many hours they sleep a night, $y$.

\begin{center}
\setlength{\tabcolsep}{7pt}\renewcommand{\arraystretch}{1.2}
\begin{tabular}{|l|c|c|c|c|c|c|c|c|}
\hline
Phone time, $x$ (hours) & 1.5 & 2 & 2.5 & 3 & 4 & 4.5 & 5.5 & 6 \\
\hline
Sleep, $y$ (hours) & 8.2 & 8.7 & 7.6 & 8.1 & 7.0 & 7.8 & 6.6 & 7.1 \\
\hline
\end{tabular}
\end{center}

\begin{enumerate}[itemsep=0.1cm]
  \item Find the equation of the regression line of $y$ on $x$, in the
  form $y = ax + b$.
  \item Interpret the meaning of the value of $b$ in this context.
  \item Find the coordinates of the mean point, $(\bar{x}, \bar{y})$.
  \item Find the equation of the regression line of $x$ on $y$.
  \item Another student, Lucia, sleeps $7.0$ hours a night. Use the
  appropriate regression line to estimate her daily phone time.
\end{enumerate}

\begin{tikzpicture}
  \draw (-0.6,0) rectangle (14.6,13.8);
  \foreach \y in {1,2,...,13} \draw[dotted] (0,\y)--(14.1,\y);
\end{tikzpicture}

\end{enumerate}
\end{document}
